% =============================================================================
% APPENDICE: Reference del GM
% =============================================================================

\chapter{Appendice: Reference del GM}
\label{ch:appendice}

\begin{multicols}{2}

\section{Tiri in Breve}
\label{sec:ref-tiri}

\begin{center}
\rowcolors{2}{LightBG}{white}
\begin{tabular}{C{1.2cm} p{5cm}}
  \toprule
  \textbf{Dado} & \textbf{Risultato} \\
  \midrule
  \textbf{6}   & Successo pieno. Effetto completo. \\
  \textbf{4–5} & Successo parziale. Complicazione. \\
  \textbf{1–3} & Fallimento. Conseguenza grave. \\
  \textbf{66}  & Critico. Effetto potenziato. \\
  \bottomrule
\end{tabular}
\end{center}

\subsection{Posizioni}
\begin{description}[leftmargin=3cm, style=nextline]
  \item[\idx{Controllata}]  Vantaggio. Conseguenze leggere.
  \item[\idx{Rischiosa}]    Standard. Conseguenze moderate.
  \item[\idx{Disperata}]    1d fisso. Conseguenze gravi.
\end{description}

\columnbreak

\section{Attributi FitD}
\label{sec:ref-attributi}

\begin{description}[leftmargin=2.5cm]
  \item[\textbf{Insight}]
    \textit{Caccia, Studio, Osserva, Ingegno}
  \item[\textbf{Prowess}]
    \textit{Finezza, Ombra, Combattimento, Forza Bruta}
  \item[\textbf{Resolve}]
    \textit{Percezione, Comando, Relazioni, Persuasione}
\end{description}

\medskip

\textbf{Resistenza:} Tira il pool dell'attributo rilevante.
\begin{itemize}
  \item Con \textbf{6}: 0 stress \ \ Con \textbf{4–5}: 1 stress \ \ Con \textbf{1–3}: 2 stress
\end{itemize}

\end{multicols}

\section{Clock Tracker — Reference Rapido}
\label{sec:ref-clocks}

\clocktrack[DoriaBlue]{DORIA — La Flotta del Tradimento [3 fasi]}{0}{3}
\clocktrack[SpinolaCrimson]{SPINOLA — Il Prezzo del Silenzio [3 fasi]}{0}{3}
\clocktrack[FieschiGreen]{FIESCHI — Le Campane di San Lorenzo [3 fasi]}{0}{3}
\clocktrack[GrimaldiPurple]{GRIMALDI — Il Prezzo dell'Autonomia [3 fasi]}{0}{3}
\clocktrack[SharedBrown]{COLPO GLOBALE — Vespri Genovesi [serve 3+ successi su 4]}{0}{4}
\clocktrack[GoldRPG]{CONTROLLO DI GENOVA [clock condiviso — 8 segmenti]}{0}{8}

\section{Ordine del Colpo Globale}
\label{sec:ref-globale}

\begin{enumerate}
  \item \textbf{Attivazione:} Il GM coordinatore annuncia l'attivazione quando
        il clock raggiunge 4 segmenti.
  \item \textbf{Dichiarazione (5 minuti):} Ogni tavolo dichiara:
        \textbf{Collaboro} / \textbf{Interferisco}.
  \item \textbf{Esecuzione:} Tutti i tavoli tirano simultaneamente
        (o con 2 minuti di offset). I risultati vengono comunicati
        al GM coordinatore.
  \item \textbf{Risoluzione:} Il GM coordinatore conta i successi
        e aggiorna il clock condiviso.
  \item \textbf{Conseguenze locali:} Ogni GM di tavolo descrive
        le conseguenze narrative nel proprio sestiere.
\end{enumerate}

\section{Ruoli nel Colpo Globale per Famiglia}
\label{sec:ref-ruoli}

\begin{center}
\rowcolors{2}{LightBG}{white}
\begin{tabularx}{\textwidth}{L{2.2cm} L{3cm} X L{3.2cm}}
  \toprule
  \textbf{Famiglia} & \textbf{Zona} & \textbf{Compito} & \textbf{Tiro} \\
  \midrule
  Doria    & Prè e porto vecchio    & Bloccare le truppe francesi dal Castelletto & Forza + Soldati \\
  Spinola  & Rete di comunicazione  & Far circolare il segnale concordato tra i sestieri & Intrigo + Influenza \\
  Fieschi  & Quartieri popolari     & Mobilitare i cittadini perché escano di casa & Influenza + Reputazione \\
  Grimaldi & Logistica nascosta     & Garantire il flusso di armi e comunicazioni & Commercio + Influenza \\
  \bottomrule
\end{tabularx}
\end{center}

\section{Vulnerabilità Incrociate}
\label{sec:ref-vulnerabilita}

\begin{center}
\rowcolors{2}{LightBG}{white}
\begin{tabular}{L{2.8cm} L{2.8cm} p{8cm}}
  \toprule
  \textbf{Colpo} & \textbf{Chi può sabotarlo} & \textbf{Come} \\
  \midrule
  Doria    & Grimaldi  & Controllano le rotte: la Fase~1 (ricognizione) fallisce automaticamente. \\
  Spinola  & Fieschi   & Canali ecclesiastici vicino a Trivulzio: le prove vengono «smarrite». \\
  Fieschi  & Doria     & Controllano i sestieri vicini: distraggono il vescovo con crisi portuale. \\
  Grimaldi & Spinola   & Agente vicino al rappresentante francese: intercettano il messaggio alla Fase~1. \\
  \bottomrule
\end{tabular}
\end{center}

\section{Risorse Iniziali a Confronto}
\label{sec:ref-risorse}

\begin{center}
\rowcolors{2}{LightBG}{white}
\begin{tabular}{L{2.5cm} C{1.5cm} C{1.5cm} C{1.5cm} C{1.5cm}}
  \toprule
  \textbf{Famiglia} & \textbf{Moneta} & \textbf{Rep.} & \textbf{Influ.} & \textbf{Sold.} \\
  \midrule
  Doria    & 4 & 5 & 3 & 4 \\
  Spinola  & 6 & 4 & 5 & 2 \\
  Fieschi  & 2 & 5 & 6 & 3 \\
  Grimaldi & 5 & 3 & 4 & 2 \\
  \bottomrule
\end{tabular}
\end{center}

\section{Note per il GM Coordinatore}
\label{sec:ref-coordinatore}

\begin{rulebox}[GoldRPG]{Comunicazione tra tavoli}
Il GM coordinatore gestisce il clock condiviso e coordina le fasi del Colpo Globale.
Strumenti consigliati per la comunicazione in tempo reale:
\textbf{Discord} (canale \texttt{\#clock-condiviso}) oppure
\textbf{Telegram} (gruppo GM only).

Aggiornare il clock visivamente dopo ogni evento significativo:
usare un bot o un messaggio pinnato che tutti i GM possono vedere in tempo reale.
\end{rulebox}

\begin{rulebox}[GoldRPG]{Segreti dei PG}
I segreti nelle schede personaggio sono evidenziati in rosso.
Si consiglia di distribuire le schede \textbf{separatamente} per giocatore
\textit{prima} della sessione, con istruzione esplicita di non condividere
la sezione rossa con i compagni di tavolo.

I segreti sono calibrati per creare tensione \textit{dentro} il tavolo
oltre che tra tavoli. Non è necessario che emergano tutti — anche solo
uno o due segreti rivelati in gioco arricchiscono la narrazione.
\end{rulebox}

\begin{rulebox}[GoldRPG]{Calibrazione della difficoltà}
Il clock a 8 segmenti è calibrato per 4 tavoli attivi per \textbf{circa 3 ore}.
La struttura temporale consigliata è:
\begin{itemize}
  \item \textbf{0:00 – 0:45} \ Fase 1–2: apertura, primo Colpo Riservato
  \item \textbf{0:45 – 1:30} \ Fase 3–4: tensione, attivazione Colpo Globale
  \item \textbf{1:30 – 2:15} \ Fase 5–6: escalation, interferenze
  \item \textbf{2:15 – 3:00} \ Fase 7–8: climax e epilogo
\end{itemize}
Con 3 tavoli, ridurre a 6 segmenti.
Con 2 tavoli, ridurre a 4 segmenti e saltare il Colpo Globale.
Il Colpo Globale si attiva al 50\% del clock (segmento 4 su 8, o 3 su 6).
\end{rulebox}
